%! Unit Opener: Functions as Models — Unit 2, Lesson 2.0 — Group Activity
\documentclass[10pt]{article}
\usepackage{atda-article}
\usepackage{atda-boxes}
\newcommand{\TallMath}[1]{$\displaystyle #1\rule[-1.4em]{0pt}{3.2em}$}
\setlength{\workrowsep}{6pt}

\begin{document}

\pageheader{Unit 2, Lesson 2.0}{Group Activity --- Reading the Day off a Graph}
% NAMESTRIP: no \namedateperiod here. The student writes their name once, on the
% cover this component is stapled behind. See references/conventions.md.

\vspace{0.15in}

\begin{scenariobox}[The Scenario]{royal}
\small
Today nobody hands you an equation. You get pictures, and every answer is in the picture. Work
with your group. \textbf{Everyone starts at Tier R.} When your whole group can do Tier R without
help, move on.
\end{scenariobox}

\vspace{0.10in}

\boxguard[30]
\begin{tcolorbox}[colback=white, colframe=black!40, title=\textbf{Tier R --- Remediate},
    fonttitle=\bfseries, arc=2mm, left=3mm, right=3mm, top=2mm, bottom=2mm, breakable]
\small
A rideshare charges a flat pickup fee plus a fixed amount per mile. Its cost $C$, in dollars, for
a ride of $m$ miles is graphed below.

\begin{center}
\begin{tikzpicture}
\begin{axis}[
    width=9.6cm, height=4.8cm,
    xlabel={$m$ (miles)}, ylabel={$C$ (dollars)},
    xmin=0, xmax=10.5, ymin=0, ymax=25,
    xtick={0,2,4,6,8,10}, ytick={0,5,10,15,20},
    grid=both, grid style={linegray!45}, axis lines=left,
    tick label style={font=\scriptsize}, label style={font=\scriptsize}]
  \addplot[royal, very thick, domain=0:10, samples=2]{2*x+3};
\end{axis}
\end{tikzpicture}
\end{center}

\begin{enumerate}[leftmargin=1.6em, itemsep=4pt, topsep=2pt]
  \item Read $C(4)$ off the graph. What does that number mean for a rider?

\writelines{1}

  \item Where does the graph meet the vertical axis, and what does that value mean before the
        car has gone anywhere?

\writelines{1}

  \item A rider is charged \$13. How many miles was the ride?
\begin{work}
  2m+3 &= 13 \\
    2m &= 10 \\
     m &= 5
\end{work}
\end{enumerate}
\end{tcolorbox}

\vspace{0.10in}

\boxguard[30]
\begin{tcolorbox}[colback=white, colframe=black!40,
    title=\textbf{Tier A --- Approaching Proficiency},
    fonttitle=\bfseries, arc=2mm, left=3mm, right=3mm, top=2mm, bottom=2mm, breakable]
\small
$B$ is the charge left in a phone's battery, in percent, $t$ hours after $7$ a.m.

\begin{center}
\begin{tikzpicture}
\begin{axis}[
    width=11.6cm, height=5.2cm,
    xlabel={$t$ (hours after 7 a.m.)}, ylabel={$B$ (percent charge)},
    xmin=0, xmax=10.5, ymin=0, ymax=110,
    xtick={0,1,2,3,4,5,6,7,8,9,10}, ytick={0,20,40,60,80,100},
    grid=both, grid style={linegray!45}, axis lines=left,
    tick label style={font=\scriptsize}, label style={font=\scriptsize}]
  \addplot[royal, very thick, mark=*, mark size=1.6pt]
    coordinates {(0,100) (2,70) (3,70) (7,30) (9,80) (10,65)};
\end{axis}
\end{tikzpicture}
\end{center}

\begin{enumerate}[leftmargin=1.6em, itemsep=4pt, topsep=2pt]
  \item State the domain and the range of this function, with units.

\writelines{1}

  \item List every interval on which $B$ is increasing, decreasing, and constant.

\writelines{2}

  \item Give the absolute maximum and the absolute minimum --- the value \emph{and} where each
        one happens. Then convert each $t$ to a clock time.

\writelines{2}

  \item Between $t=7$ and $t=9$ the battery went from $30\%$ to $80\%$. What was happening, and
        at what rate, in percent per hour?
\begin{work}
  \text{rate} &= \dfrac{80-30}{9-7} \\
              &= 25 \text{ percent per hour}
\end{work}

  \item The phone was in airplane mode for one stretch of the morning. Which piece of the graph
        shows it, and what is that piece called?

\writelines{1}
\end{enumerate}
\end{tcolorbox}

\vspace{0.10in}

\boxguard[30]
\begin{tcolorbox}[colback=white, colframe=black!40, title=\textbf{Tier E --- Extension},
    fonttitle=\bfseries, arc=2mm, left=3mm, right=3mm, top=2mm, bottom=2mm, breakable]
\small

\begin{enumerate}[leftmargin=1.6em, itemsep=4pt, topsep=2pt]
  \item \textbf{Critique the work.} Two students answered questions about the battery graph.
        Each is wrong. Name the error and give the correct answer.
        \begin{enumerate}[label=(\alph*), leftmargin=1.6em, itemsep=3pt, topsep=3pt]
          \item ``Domain: $30 \le B \le 100$. \quad Range: $0 \le t \le 10$.''

\writelines{1}

          \item ``The absolute maximum is $80\%$, because that is the highest the battery ever
                charged back up to.''

\writelines{1}
        \end{enumerate}

  \item \textbf{Which family?} Each table below comes from one of the three families. Decide
        which, and give the evidence you used.
        \smallskip

        \renewcommand{\arraystretch}{1.25}
        \begin{tabular}{@{}|c|c|c|c|c|@{}}
        \hline
        $x$ & $0$ & $1$ & $2$ & $3$ \\ \hline
        $A$ & $5$ & $8$ & $11$ & $14$ \\ \hline
        $B$ & $4$ & $8$ & $16$ & $32$ \\ \hline
        $C$ & $0$ & $3$ & $12$ & $27$ \\ \hline
        \end{tabular}
        \smallskip

\writelines{3}

  \item \textbf{Justify.} A classmate says, ``Any graph that keeps going up is exponential.'' Use
        two of the tables above to show that is wrong.

\writelines{2}
\end{enumerate}
\end{tcolorbox}

\end{document}
